% Part: proof-theory
% Chapter: normalization
% Section: introduction

\documentclass[../../../include/open-logic-section]{subfiles}

\begin{document}

\olfileid{pt}{nor}{int}

\olsection{પરિચય}

જો તમે અનુસરણ~$!A \lif !B$ને
!!{prove}d કર્યું હોય, તો
\Elim{\lif} વાપરીને
તેને~$!B$ની !!a{proof}માં
કામે લગાડી શકો.
અલબત્ત, આ માટે~$!A$ની
!!a{proof}~$\delta_1$ પણ
જોઈએ. અનુસરણ~$!A \lif !B$ની
તમારી !!{proof}~$\delta$
મોટે ભાગે~\Intro{\lif}થી
પૂરી થતી હશે:
તે નિયમ ખુલ્લી
ધારણા~$!A$ હેઠળની
$!B$ની !!a{proof}~$\delta_2$ને
$!A \lif !B$ની સાબિતીમાં
ફેરવે છે. એમ હોય તો
અંતિમ !!{proof}નું
સ્વરૂપ આ છે:
\[
\AxiomC{$\Discharge{!A}{x}$}
\RightLabel{$\delta_2$}
\DeduceC{$!B$}
\DischargeRule{\Intro{\lif}}{x}
\UnaryInfC{$!A \lif !B$}
\AxiomC{}
\RightLabel{$\delta_1$}
\DeduceC{$!A$}
\RightLabel{\Elim{\lif}}
\BinaryInfC{$!B$}
\DisplayProof
\]
\sourcecorrection{OLMD-252}{મૂળ ગદ્યમાં $\delta_1$ને પહેલે $!A$ની અને પછી $!A \lif !B$ની સાબિતી કહેવામાં આવી છે; વૃક્ષના બીજા આધાર ઉપર પણ $\delta_2$ ખોટી રીતે પુનરાવર્તિત છે. અહીં અનુસરણની સાબિતી $\delta$ અને $!A$નો આધાર $\delta_1$ છે.}
પહેલેથી મળેલી
$!A \lif !B$ની !!a{proof}
ફરી વાપરવાથી તમારી
રીત કાર્યક્ષમ છે,
પરંતુ તેની !!{proof}
કાર્યક્ષમ નથી.
$!A \lif !B$ને રજૂ
કરીને તરત જ તેનો
લોપ કરવો ચકરાવો લાગે
છે. $!B$ને !!{prove}
કરવાની વધુ સીધી
રીત હોવી જોઈએ.

ખરેખર, એવી રીત હંમેશાં
હોય છે. ઉપરના જેવા
``ચકરાવા''—!!a{introduction}
નિયમ પછી તરત
!!a{elimination} નિયમ—સ્વાભાવિક
નિગમનની !!{proof}sમાંથી
હંમેશાં દૂર કરી શકાય.
તેને \emph{સામાન્યીકરણ}
કહે છે અને મળતી
સાબિતીને \emph{સામાન્ય}
!!{proof} (અથવા
\emph{સામાન્ય સ્વરૂપમાં}
હોય એવી !!a{proof})
કહે છે.

આ પરિણામની સાબિતી,
સિક્વન્ટ-કલનના કટ-નિવારણ
પ્રમેયની જેમ, થોડી
જટિલ છે અને તેમાં
ઘણા કિસ્સા જોવા પડે છે.
અંતઃપ્રજ્ઞાવાદી~\Log{N1i}
કરતાં શાસ્ત્રીય~\Log{N1c}
માટે તે સાબિત કરવું
વધુ મુશ્કેલ છે.
સિક્વન્ટ-કલનમાં \Cut{}
નિયમ વાપરીને તેની
વિના કરતાં વધુ
પરિણામો !!{prove}
કરી શકાતાં નથી એવું
કટ-નિવારણ બતાવે છે;
તેવી જ રીતે,
ચકરાવા વાપરીને તેમને
ટાળવા કરતાં વધુ
પરિણામો સાબિત કરી
શકાતાં નથી એવું
સામાન્યીકરણ બતાવે છે.
બીજી સામ્યતાઓ પણ છે:
એક જ પરિણામની ચકરાવાવાળી
સૌથી ટૂંકી !!{proof}
કરતાં સામાન્ય
!!{proof} ઘણી લાંબી
હોઈ શકે છે; તેવી જ
રીતે, સિક્વન્ટ-કલનમાં
કટવાળી !!{proof}s
સૌથી ટૂંકી કટ-રહિત
!!{proof} કરતાં ઘણી
ટૂંકી હોઈ શકે છે.
ઉપ-!!{formula}
ગુણધર્મ કહે છે કે
પરિણામના ઉપ-!!{formula}s
વાપરીને જ તે પરિણામ
!!{prove} કરતી
!!{proof} હંમેશાં
મેળવી શકાય છે.
સિક્વન્ટ-કલનમાં આ
ગુણધર્મ કટ-નિવારણમાંથી
અને સ્વાભાવિક નિગમનમાં
સામાન્યીકરણમાંથી ફલિત
થાય છે.

\end{document}
